\section{破产法基本制度}
\begin{enumerate}
  \item 主观题。
  \item \zd{公平受偿}。避免个别的清偿。
  \item 申请人（债权人、债务人）\yjt
        法院受理 \yjt
        指定管理人 \yjt
        通知公告 \yjt
        申报债权 \yjt
        第一次债权人会议 \yjt
        宣告破产 \yjt
        拟定财产的变价方案。
  \item 实践清偿比例 30\% 就很不错了。
  \item 破产原因。受理标准；宣告破产标准。
        \begin{enumerate}
          \item 不能清偿到期债务 + 资不抵债。形式判断。
          \item 不能清偿到期债务 + 明显缺乏清偿能力。实质判断。
        \end{enumerate}
  \item 申请。
        \begin{enumerate}
          \item 主体。
                \begin{enumerate}
                  \item 债权人只要求不能清偿到期债务。只能申请清算、重整。
                  \item 债务人申请。破产原因。金蝉脱壳。
                  \item 清算组转过来。
                  \item 执行转破产。主观题。
                \end{enumerate}
          \item 审查。债务人异议期；是否主管管辖；破产原因；到期等初步审查。
          \item 受理。很少考。
          \item 撤回、驳回。太细节了。用的都是裁定。
        \end{enumerate}
  \item 受理。是破产开始的标志。
        \begin{enumerate}
          \item 管理人接管事务。
                \begin{enumerate}
                  \item 向管理人（不是董监高）清偿债务。
                  \item 有双方未履行（完毕）的双务合同。
                        \begin{enumerate}
                          \item 管理人决定继续 \hz 解除。
                          \item 受理 2 月 \hz 催告日 30 天后默认解除。
                          \item 重大影响报法院、债权人。
                        \end{enumerate}
                \end{enumerate}
          \item 保障公平受偿权。
                \begin{enumerate}
                  \item 受理后禁止个别清偿。无效。
                  \item 保全措施清除，执行中止，银行停止计息，滞纳金。
                \end{enumerate}
          \item 对诉讼影响。
                \begin{enumerate}
                  \item 受理前民事诉讼中止。
                  \item 受理后不接受清偿诉讼，只接受确认债权额诉讼。
                  \item 走集中管辖。处理破产的法院管辖。
                \end{enumerate}
        \end{enumerate}
  \item 破产管理人。\ydy 董事会、高管。
        \begin{enumerate}
          \item 法院指定，辞职法院许可，决定是否更换。法院居主导。
          \item 资格。清算组、中介机构、个人（律师、会计师）。
          \item 职责。财产、报告、方案。可以继续经营。管理人决定 \yjt 债权人会议。
        \end{enumerate}
  \item 债权会议。\ydy 股东会。
        \begin{enumerate}
          \item 一般事项。出席会议的人数过半数，债权金额占无财产担保债权总额的 1/2。
          \item 特殊。重整、和解。1/2 + 2/3。
        \end{enumerate}
  \item 债权人委员会。监督机构。
  \item \zd{有重大影响的财产处分制度}。
        \begin{enumerate}
          \item 涉及不动产。
          \item 探矿权、采矿权、知识产权等财产权转让。
          \item 全部库存 \hz 营业的转让。整体转让。
          \item 借款。
          \item 继续履行。
        \end{enumerate}
\end{enumerate}
